\documentclass[10pt,a4paper]{amsart}
\usepackage{amsmath,amssymb,amsthm,graphicx}
\usepackage{verbatim}
\newenvironment{reference}{\comment}{\endcomment}
%\newenvironment{reference}{}{}
\newenvironment{slogan}{\comment}{\endcomment}
\newenvironment{history}{\comment}{\endcomment}

% For commutative diagrams we use Xy-pic
\usepackage[all]{xy}

% We use 2cell for 2-commutative diagrams.
\xyoption{2cell}
% Theorem environments.
%
\theoremstyle{plain}
\newtheorem{theorem}[subsection]{Theorem}
\newtheorem{proposition}[subsection]{Proposition}
\newtheorem{lemma}[subsection]{Lemma}

\theoremstyle{definition}
\newtheorem{definition}[subsection]{Definition}
\newtheorem{example}[subsection]{Example}
\newtheorem{exercise}[subsection]{Exercise}
\newtheorem{situation}[subsection]{Situation}

\theoremstyle{remark}
\newtheorem{remark}[subsection]{Remark}
\newtheorem{remarks}[subsection]{Remarks}

\numberwithin{equation}{subsection}

% Macros
%
\def\lim{\mathop{\mathrm{lim}}\nolimits}
\def\colim{\mathop{\mathrm{colim}}\nolimits}
\def\Spec{\mathop{\mathrm{Spec}}}
\def\Hom{\mathop{\mathrm{Hom}}\nolimits}
\def\Ext{\mathop{\mathrm{Ext}}\nolimits}
\def\SheafHom{\mathop{\mathcal{H}\!\mathit{om}}\nolimits}
\def\SheafExt{\mathop{\mathcal{E}\!\mathit{xt}}\nolimits}
\def\Sch{\mathit{Sch}}
\def\Mor{\mathop{\mathrm{Mor}}\nolimits}
\def\Ob{\mathop{\mathrm{Ob}}\nolimits}
\def\Sh{\mathop{\mathit{Sh}}\nolimits}
\def\NL{\mathop{N\!L}\nolimits}
\def\CH{\mathop{\mathrm{CH}}\nolimits}
\def\proetale{{pro\text{-}\acute{e}tale}}
\def\etale{{\acute{e}tale}}
\def\QCoh{\mathit{QCoh}}
\def\Ker{\mathop{\mathrm{Ker}}}
\def\Im{\mathop{\mathrm{Im}}}
\def\Coker{\mathop{\mathrm{Coker}}}
\def\Coim{\mathop{\mathrm{Coim}}}

% Boxtimes
%
\DeclareMathSymbol{\boxtimes}{\mathbin}{AMSa}{"02}

%
% Macros for moduli stacks/spaces
%
\def\QCohstack{\mathcal{QC}\!\mathit{oh}}
\def\Cohstack{\mathcal{C}\!\mathit{oh}}
\def\Spacesstack{\mathcal{S}\!\mathit{paces}}
\def\Quotfunctor{\mathrm{Quot}}
\def\Hilbfunctor{\mathrm{Hilb}}
\def\Curvesstack{\mathcal{C}\!\mathit{urves}}
\def\Polarizedstack{\mathcal{P}\!\mathit{olarized}}
\def\Complexesstack{\mathcal{C}\!\mathit{omplexes}}
% \Pic is the operator that assigns to X its picard group, usage \Pic(X)
% \Picardstack_{X/B} denotes the Picard stack of X over B
% \Picardfunctor_{X/B} denotes the Picard functor of X over B
\def\Pic{\mathop{\mathrm{Pic}}\nolimits}
\def\Picardstack{\mathcal{P}\!\mathit{ic}}
\def\Picardfunctor{\mathrm{Pic}}
\def\Deformationcategory{\mathcal{D}\!\mathit{ef}}
\usepackage{cleveref}
\crefname{originallemma}{original lemma}{original lemmas}
\Crefname{originallemma}{Original lemma}{Original lemmas}
\makeatletter
\newlabel{derived-lemma-representable-homological}{{0149}{source}{}{sourcepointer-1}{}}
\newlabel{derived-lemma-representable-homological@cref}{{[originallemma][0][]0149}{[1][0][]source}{}{}{}}
\makeatother
\crefname{originallemma}{original lemma}{original lemmas}
\Crefname{originallemma}{Original lemma}{Original lemmas}
\makeatletter
\newlabel{algebra-lemma-directed-colimit-exact}{{00DB}{source}{}{sourcepointer-2}{}}
\newlabel{algebra-lemma-directed-colimit-exact@cref}{{[originallemma][0][]00DB}{[1][0][]source}{}{}{}}
\makeatother
\crefname{originalreference}{original reference}{original references}
\Crefname{originalreference}{Original reference}{Original references}
\makeatletter
\newlabel{derived-remark-truncation-distinguished-triangle}{{08J5}{source}{}{sourcepointer-3}{}}
\newlabel{derived-remark-truncation-distinguished-triangle@cref}{{[originalreference][0][]08J5}{[1][0][]source}{}{}{}}
\makeatother
\crefname{originallemma}{original lemma}{original lemmas}
\Crefname{originallemma}{Original lemma}{Original lemmas}
\makeatletter
\newlabel{coherent-lemma-homs-over-open}{{01YB}{source}{}{sourcepointer-4}{}}
\newlabel{coherent-lemma-homs-over-open@cref}{{[originallemma][0][]01YB}{[1][0][]source}{}{}{}}
\makeatother
\crefname{originalproposition}{original proposition}{original propositions}
\Crefname{originalproposition}{Original proposition}{Original propositions}
\makeatletter
\newlabel{perfect-proposition-DCoh}{{0FDB}{source}{}{sourcepointer-5}{}}
\newlabel{perfect-proposition-DCoh@cref}{{[originalproposition][0][]0FDB}{[1][0][]source}{}{}{}}
\makeatother
\crefname{originallemma}{original lemma}{original lemmas}
\Crefname{originallemma}{Original lemma}{Original lemmas}
\makeatletter
\newlabel{coherent-lemma-Artin-Rees}{{01YA}{source}{}{sourcepointer-6}{}}
\newlabel{coherent-lemma-Artin-Rees@cref}{{[originallemma][0][]01YA}{[1][0][]source}{}{}{}}
\makeatother
\crefname{originallemma}{original lemma}{original lemmas}
\Crefname{originallemma}{Original lemma}{Original lemmas}
\makeatletter
\newlabel{properties-lemma-extend-finite-presentation}{{01PI}{source}{}{sourcepointer-7}{}}
\newlabel{properties-lemma-extend-finite-presentation@cref}{{[originallemma][0][]01PI}{[1][0][]source}{}{}{}}
\makeatother
\crefname{originallemma}{original lemma}{original lemmas}
\Crefname{originallemma}{Original lemma}{Original lemmas}
\makeatletter
\newlabel{coherent-lemma-power-ideal-kills-sheaf}{{01Y9}{source}{}{sourcepointer-8}{}}
\newlabel{coherent-lemma-power-ideal-kills-sheaf@cref}{{[originallemma][0][]01Y9}{[1][0][]source}{}{}{}}
\makeatother
\title{The colimit of Hom out of a Deligne system computes Hom on the open}
\author{The Stacks Project Authors}
\date{stacks-project:duality.tex:local-snapshot-20261012; GNU FDL 1.2 or later}
\begin{document}
\maketitle
\noindent\textit{Editorial note.} All delivered mathematical text is the author's own, copied from the pinned Stacks Project chapter file duality.tex; no statement or proof step was rewritten, reordered or removed.\par\medskip
\noindent\textit{Editorial note.} Preamble provenance: the pinned chapter file opens with an input of the shared preamble and carries no preamble text of its own, so the project's single Stacks-wide preamble staged at sources/stacks-duality/preamble.tex is excerpted verbatim under the recipe fixed by LANE\_BRIEF 22.1, as three spans, lines 8-13 (verbatim package and the reference/slogan/history environments), lines 14-18 (Xy-pic setup) and lines 154-217 (the native theorem declarations and mathematical macros). Nothing was added, removed or redefined; the wrapper class and the amsmath/amssymb/amsthm/graphicx base packages are supplied by the assembler.\par\medskip
\noindent\textit{Editorial note.} The counted statement and proof contain no cross reference, no citation and no included file, so the delivered body is self contained.\par\medskip
\noindent\textit{Editorial note.} Canonical location of the counted proposition: Stacks tag 0G2H (sources/stacks-duality/pinned-original-source.tex line 7920).\par\medskip
\noindent\textit{Allowed standard prerequisites (editorial).}\begin{itemize}
\item Deligne systems and their constant values on the open. Same chapter, the definition of a Deligne system preceding this lemma, invoked in the statement. Fixes the compatible system (K\_n) and its value K.
\item Exactness of filtered colimits and the long exact sequence of Hom. Derived Categories, Lemma labelled derived-lemma-representable-homological and Algebra, Lemma labelled algebra-lemma-directed-colimit-exact, cited at proof lines 7939-7943. Justifies the 5-lemma reduction along distinguished triangles.
\item Truncation triangles and the identification of the bounded coherent derived category with the derived category of coherent sheaves. Derived Categories, Remark labelled derived-remark-truncation-distinguished-triangle, cited at proof lines 7944-7946, and Derived Categories of Schemes, Proposition labelled perfect-proposition-DCoh, cited at proof lines 7998-8002. Reduces to one nonzero cohomology sheaf and lets Ext be computed in the abelian category.
\item Artin-Rees and the vanishing of Hom over the open. Cohomology of Schemes, Lemmas labelled coherent-lemma-Artin-Rees, coherent-lemma-homs-over-open and coherent-lemma-power-ideal-kills-sheaf, cited at proof lines 8016-8018, 7968-7970 and 8078-8080. Supplies the pro-isomorphism of the two inverse systems, the base case i = 0 and the final containment.
\item Extending a coherent module from an open. Properties, Lemma labelled properties-lemma-extend-finite-presentation, cited at proof lines 8034-8038. Realises the auxiliary map over U by a map over X in the surjectivity argument.
\end{itemize}\medskip
\ifdefined\hypertarget\hypertarget{sourcepointer-1}{}\fi
\noindent\textit{Source reference.} Original lemma 0149: sources/stacks-chow/tags.txt, source lines 1464--1464. Cross-chapter standard prerequisite located in the project tag list staged in the local snapshot; invoked at proof line 7948 of the counted proof. It is retained only as a navigation pointer; every load-bearing step is inside the delivered span.\par\medskip
\ifdefined\hypertarget\hypertarget{sourcepointer-2}{}\fi
\noindent\textit{Source reference.} Original lemma 00DB: sources/stacks-chow/tags.txt, source lines 491--491. Cross-chapter standard prerequisite located in the project tag list staged in the local snapshot; invoked at proof line 7949 of the counted proof. It is retained only as a navigation pointer; every load-bearing step is inside the delivered span.\par\medskip
\ifdefined\hypertarget\hypertarget{sourcepointer-3}{}\fi
\noindent\textit{Source reference.} Original reference 08J5: sources/stacks-chow/tags.txt, source lines 10838--10838. Cross-chapter standard prerequisite located in the project tag list staged in the local snapshot; invoked at proof line 7954 of the counted proof. It is retained only as a navigation pointer; every load-bearing step is inside the delivered span.\par\medskip
\ifdefined\hypertarget\hypertarget{sourcepointer-4}{}\fi
\noindent\textit{Source reference.} Original lemma 01YB: sources/stacks-chow/tags.txt, source lines 2550--2550. Cross-chapter standard prerequisite located in the project tag list staged in the local snapshot; invoked at proof line 7985 of the counted proof. It is retained only as a navigation pointer; every load-bearing step is inside the delivered span.\par\medskip
\ifdefined\hypertarget\hypertarget{sourcepointer-5}{}\fi
\noindent\textit{Source reference.} Original proposition 0FDB: sources/stacks-chow/tags.txt, source lines 19211--19211. Cross-chapter standard prerequisite located in the project tag list staged in the local snapshot; invoked at proof line 7997 of the counted proof. It is retained only as a navigation pointer; every load-bearing step is inside the delivered span.\par\medskip
\ifdefined\hypertarget\hypertarget{sourcepointer-6}{}\fi
\noindent\textit{Source reference.} Original lemma 01YA: sources/stacks-chow/tags.txt, source lines 2549--2549. Cross-chapter standard prerequisite located in the project tag list staged in the local snapshot; invoked at proof line 8016 of the counted proof. It is retained only as a navigation pointer; every load-bearing step is inside the delivered span.\par\medskip
\ifdefined\hypertarget\hypertarget{sourcepointer-7}{}\fi
\noindent\textit{Source reference.} Original lemma 01PI: sources/stacks-chow/tags.txt, source lines 2231--2231. Cross-chapter standard prerequisite located in the project tag list staged in the local snapshot; invoked at proof line 8055 of the counted proof. It is retained only as a navigation pointer; every load-bearing step is inside the delivered span.\par\medskip
\ifdefined\hypertarget\hypertarget{sourcepointer-8}{}\fi
\noindent\textit{Source reference.} Original lemma 01Y9: sources/stacks-chow/tags.txt, source lines 2548--2548. Cross-chapter standard prerequisite located in the project tag list staged in the local snapshot; invoked at proof line 8083 of the counted proof. It is retained only as a navigation pointer; every load-bearing step is inside the delivered span.\par\medskip
\begin{lemma}
\label{lemma-lift-map}
Let $j : U \to X$ be an open immersion of Noetherian schemes.
Let $(K_n)$ be a Deligne system and denote
$K \in D^b_{\textit{Coh}}(\mathcal{O}_U)$ the value
of the constant system $(K_n|_U)$. Let $L$ be an object of
$D^b_{\textit{Coh}}(\mathcal{O}_X)$.
Then $\colim \Hom_X(K_n, L) = \Hom_U(K, L|_U)$.
\end{lemma}

\begin{proof}
Let $L \to M \to N \to L[1]$ be a distinguished triangle in
$D^b_{\textit{Coh}}(\mathcal{O}_X)$. Then we obtain
a commutative diagram
$$
\xymatrix{
\ldots \ar[r] &
\colim \Hom_X(K_n, L) \ar[r] \ar[d] &
\colim \Hom_X(K_n, M) \ar[r] \ar[d] &
\colim \Hom_X(K_n, N) \ar[r] \ar[d] &
\ldots \\
\ldots \ar[r] &
\Hom_U(K, L|_U) \ar[r] &
\Hom_U(K, M|_U) \ar[r] &
\Hom_U(K, N|_U) \ar[r] &
\ldots
}
$$
whose rows are exact by Derived Categories, Lemma
\ref{derived-lemma-representable-homological} and
Algebra, Lemma \ref{algebra-lemma-directed-colimit-exact}.
Hence if the statement of the lemma holds for
$N[-1]$, $L$, $N$, and $L[1]$ then it holds for $M$ by the 5-lemma.
Thus, using the distinguished triangles
for the canonical truncations of $L$ (see Derived Categories, Remark
\ref{derived-remark-truncation-distinguished-triangle})
we reduce to the case that $L$ has only one nonzero cohomology sheaf.

\medskip\noindent
Choose a bounded complex $\mathcal{F}^\bullet$ of coherent
$\mathcal{O}_X$-modules and a quasi-coherent ideal
$\mathcal{I} \subset \mathcal{O}_X$ cutting out $X \setminus U$
such that $K_n$ is represented by $\mathcal{I}^n\mathcal{F}^\bullet$.
Using ``stupid'' truncations we obtain compatible termwise split short
exact sequences of complexes
$$
0 \to \sigma_{\geq a + 1} \mathcal{I}^n\mathcal{F}^\bullet \to
\mathcal{I}^n\mathcal{F}^\bullet \to
\sigma_{\leq a} \mathcal{I}^n\mathcal{F}^\bullet \to 0
$$
which in turn correspond to compatible systems of distinguished
triangles in $D^b_{\textit{Coh}}(\mathcal{O}_X)$.
Arguing as above we reduce to the case where $\mathcal{F}^\bullet$
has only one nonzero term. This reduces us to the case discussed in
the next paragraph.

\medskip\noindent
Given a coherent $\mathcal{O}_X$-module $\mathcal{F}$ and a coherent
$\mathcal{O}_X$-module $\mathcal{G}$ we
have to show that the canonical map
$$
\colim \Ext^i_X(\mathcal{I}^n\mathcal{F}, \mathcal{G})
\longrightarrow
\Ext^i_U(\mathcal{F}|_U, \mathcal{G}|_U)
$$
is an isomorphism for all $i \geq 0$. For $i = 0$ this is
Cohomology of Schemes, Lemma \ref{coherent-lemma-homs-over-open}.
Assume $i > 0$.

\medskip\noindent
Injectivity. Let $\xi \in \Ext^i_X(\mathcal{I}^n\mathcal{F}, \mathcal{G})$
be an element whose restriction to $U$ is zero. We have to show there exists
an $m \geq n$ such that the restriction of $\xi$ to
$\mathcal{I}^m\mathcal{F} = \mathcal{I}^{m - n}\mathcal{I}^n\mathcal{F}$
is zero. After replacing $\mathcal{F}$ by $\mathcal{I}^n\mathcal{F}$
we may assume $n = 0$, i.e., we have
$\xi \in \Ext^i_X(\mathcal{F}, \mathcal{G})$
whose restriction to $U$ is zero. By
Derived Categories of Schemes, Proposition \ref{perfect-proposition-DCoh}
we have
$D^b_{\textit{Coh}}(\mathcal{O}_X) = D^b(\textit{Coh}(\mathcal{O}_X))$.
Hence we can compute the $\Ext$ group in the abelian category of
coherent $\mathcal{O}_X$-modules. This implies there exists an
surjection $\alpha : \mathcal{F}'' \to \mathcal{F}$ such that
$\xi \circ \alpha = 0$ (this is where we use that $i > 0$).
Set $\mathcal{F}' = \Ker(\alpha)$ so that we have a short exact
sequence
$$
0 \to \mathcal{F}' \to \mathcal{F}'' \to \mathcal{F} \to 0
$$
It follows that $\xi$ is the image of an element
$\xi' \in \Ext^{i - 1}_X(\mathcal{F}', \mathcal{G})$
whose restriction to $U$ is in the image of
$\Ext^{i - 1}_U(\mathcal{F}''|_U, \mathcal{G}|_U) \to
\Ext^{i - 1}_U(\mathcal{F}'|_U, \mathcal{G}|_U)$.
By Artin-Rees the inverse systems $(\mathcal{I}^n\mathcal{F}')$ and
$(\mathcal{I}^n \mathcal{F}'' \cap \mathcal{F}')$ are pro-isomorphic, see
Cohomology of Schemes, Lemma \ref{coherent-lemma-Artin-Rees}.
Since we have the compatible system of short exact sequences
$$
0 \to
\mathcal{F}' \cap \mathcal{I}^n\mathcal{F}'' \to
\mathcal{I}^n\mathcal{F}'' \to
\mathcal{I}^n\mathcal{F} \to 0
$$
we obtain a commutative diagram
$$
\xymatrix{
\colim \Ext^{i - 1}_X(\mathcal{I}^n\mathcal{F}'', \mathcal{G})
\ar[r] \ar[d] &
\colim \Ext^{i - 1}_X(\mathcal{F}' \cap \mathcal{I}^n\mathcal{F}'', \mathcal{G})
\ar[r] \ar[d] &
\colim \Ext^i_X(\mathcal{I}^n\mathcal{F}, \mathcal{G})
\ar[d] \\
\Ext^{i - 1}_U(\mathcal{F}''|_U, \mathcal{G}|_U) \ar[r] &
\Ext^{i - 1}_U(\mathcal{F}'|_U, \mathcal{G}|_U) \ar[r] &
\Ext^{i - 1}_U(\mathcal{F}|_U, \mathcal{G}|_U)
}
$$
with exact rows. By induction on $i$ and the comment on inverse systems above
we find that the left two vertical arrows are isomorphisms.
Now $\xi$ gives an element in the top right group which is the image
of $\xi'$ in the middle top group, which in turn maps to an element of
the bottom middle group coming from some element in the left bottom group.
We conclude that $\xi$ maps to zero in
$\Ext^i_X(\mathcal{I}^n\mathcal{F}, \mathcal{G})$
for some $n$ as desired.

\medskip\noindent
Surjectivity. Let $\xi \in \Ext^i_U(\mathcal{F}|_U, \mathcal{G}|_U)$.
Arguing as above using that $i > 0$ we can find an surjection
$\mathcal{H} \to \mathcal{F}|_U$ of coherent $\mathcal{O}_U$-modules
such that $\xi$ maps to zero in $\Ext^i_U(\mathcal{H}, \mathcal{G}|_U)$.
Then we can find a map $\varphi : \mathcal{F}'' \to \mathcal{F}$
of coherent $\mathcal{O}_X$-modules whose restriction to $U$ is
$\mathcal{H} \to \mathcal{F}|_U$, see
Properties, Lemma \ref{properties-lemma-extend-finite-presentation}.
Observe that the lemma doesn't guarantee $\varphi$ is surjective
but this won't matter (it is possible to pick a surjective $\varphi$
with a little bit of additional work).
Denote $\mathcal{F}' = \Ker(\varphi)$. The short exact sequence
$$
0 \to \mathcal{F}'|_U \to \mathcal{F}''|_U \to \mathcal{F}|_U \to 0
$$
shows that $\xi$ is the image of $\xi'$ in
$\Ext^{i - 1}_U(\mathcal{F}'|_U, \mathcal{G}|_U)$.
By induction on $i$ we can find an $n$ such that
$\xi'$ is the image of some $\xi'_n$ in
$\Ext^{i - 1}_X(\mathcal{I}^n\mathcal{F}', \mathcal{G})$.
By Artin-Rees we can find an $m \geq n$ such that
$\mathcal{F}' \cap \mathcal{I}^m\mathcal{F}'' \subset
\mathcal{I}^n\mathcal{F}'$. Using the short exact sequence
$$
0 \to \mathcal{F}' \cap \mathcal{I}^m\mathcal{F}'' \to
\mathcal{I}^m\mathcal{F}'' \to \mathcal{I}^m\Im(\varphi) \to 0
$$
the image of $\xi'_n$ in
$\Ext^{i - 1}_X(\mathcal{F}' \cap \mathcal{I}^m\mathcal{F}'', \mathcal{G})$
maps by the boundary map to an element $\xi_m$ of
$\Ext^i_X(\mathcal{I}^m\Im(\varphi), \mathcal{G})$ which maps
to $\xi$. Since $\Im(\varphi)$ and $\mathcal{F}$ agree over $U$
we see that $\mathcal{F}/\mathcal{I}^m\Im(\varphi)$ is supported
on $X \setminus U$. Hence there exists an $l \geq m$
such that $\mathcal{I}^l\mathcal{F} \subset \mathcal{I}^m\Im(\varphi)$, see
Cohomology of Schemes, Lemma \ref{coherent-lemma-power-ideal-kills-sheaf}.
Taking the image of $\xi_m$ in
$\Ext^i_X(\mathcal{I}^l\mathcal{F}, \mathcal{G})$ we win.
\end{proof}

\end{document}
